\begin{vidu} %[2P4-7.3B][Loại 1]
Hạt electron bay vào từ trường đều $B = 3,14. 10^{-4}$ T, theo phương vuông góc với đường sức từ. Biết tốc độ của hạt là $8. 10^6$ m/s. Khối lượng $m_e = 9,1. 10^{-31}$ kg.  Electron chuyển động theo quỹ đạo tròn có bán kính là
\choice
{1,45 cm}
{0,06 m}
{\True 14,5 cm}
{0,06 cm}
\loigiai{Ta có: $R=\dfrac{mv}{\left| q \right|B}=\dfrac{9,{{1. 10}^{-31}}{{. 8. 10}^{6}}}{1,{{6. 10}^{-19}}. 3,{{14. 10}^{-4}}}=0,145\ m=14\ cm$}
\end{vidu}
\begin{vidu} %[2P4-7.3K][Loại 3]
	Một proton chuyển động theo quỹ đạo tròn bán kính 5 cm trong một từ trường đều có $B = 0,2$ T. Biết $m_p = 1,672. 10^{-27}$ kg. Tốc độ và chu kì chuyển động của proton là
	\choice
	{$v = 4,88. 10^5$ m/s, $T = 0,656$ $\mu s$}
	{$v = 9,57. 10^4$ m/s, $T = 0,328$ $\mu s$}
	{\True $v = 9,57. 10^5$ m/s, $T = 0,328$ $\mu s$}
	{$v = 4,88. 10^4$ m/s, $T = 0,656$ $\mu s$}
	\loigiai{$v=\dfrac{\left| q \right|BR}{m} = 9,57. 10^5$ m/s $\Rightarrow T=\dfrac{2\pi R}{v}= 3,3. 10^{-7}$ s}
\end{vidu}
\begin{vidu} %[2P4-7.3B][Loại 1]
	Một electron bay với vận tốc $v = 2,4.10^6$ m/s vào trong từ trường đều $B = 1$ T theo hướng hợp với $\vec{B}$ một góc $60\degree$. Bán kính quỹ đạo chuyển động của electron là
	\choice
	{6,25 $\mu m$}
	{1,182 $\mu m$}
	{\True 15,76 $\mu m$}
	{0,625 $\mu m$}
	\loigiai{Ta có: ${F_L} = q{v_ \bot}.B = q.v.\sin \alpha .B;\ {F_L} = {F_{ht}} \to qvB\sin \alpha = \dfrac{m{v^2}}{R} \Rightarrow R = \dfrac{mv}{qB\sin \alpha} = \dfrac{{9,{{1.10}^{- 31}}.2,{{4.10}^6}}}{{1,6.10^{-19}.1.\sin {{60}^o}}} = 1,{576.10^{- 5}}m = 15,76\ \mu m$}
\end{vidu}
\begin{vidu} %[2P4-7.3B][Loại 2]
	Hai điện tích có cùng điện tích, cùng khối lượng bay vuông với các đường cảm ứng vào cùng một từ trường đều. Bỏ qua độ lớn của trọng lực. Điện tích một bay với vận tốc 1000 m/s thì có bán kính quỹ đạo 20 cm. Điện tích hai bay với vận tốc 1200 m/s thì có bán kính quỹ đạo là
	\choice
	{20 cm}
	{21 cm}
	{\True 24 cm}
	{$\dfrac{200}{11}$ cm}
	\loigiai{Ta có: $R=\dfrac{mv}{\left|q\right|B}\Rightarrow \dfrac{R_2}{R_1}=\dfrac{v_2}{v_1}\Rightarrow R_2=R_1\dfrac{v_2}{v_1}=24\ cm$}
\end{vidu}
\begin{vidu} %[2P4-7.3K][Loại 2]
	Hai điện tích $q_1 = 8\ \mu C$ và $q_2 = -2\ \mu C$ có cùng khối lượng và ban đầu chúng bay cùng hướng cùng vận tốc vào một từ trường đều. Điện tích ${q_1}$ chuyển động cùng chiều kim đồng hồ với bán kính quỹ đạo 4 cm. Điện tích ${q_2}$ chuyển động
	\choice
	{\True ngược chiều kim đồng hồ với bán kính 16 cm}
	{cùng chiều kim đồng hồ với bán kính 16 cm}
	{ngược chiều kim đồng hồ với bán kính 8 cm}
	{cùng chiều kim đồng hồ với bán kính 8 cm}
	\loigiai{\begin{itemize}
			\item Ta có:
			\begin{align*}
				R=\dfrac{mv}{\left|q\right|B}\Rightarrow \dfrac{R_2}{R_1}=\dfrac{\left|{q_1}\right|}{\left|{q_2}\right|}\Rightarrow \dfrac{R_2}{R_1}=4\Rightarrow R_2=4\text{R}_1=16\ cm.
			\end{align*}
			\item Vì $q_2<0$ nên nó sẽ chuyển động ngược chiều kim đồng hồ.
	\end{itemize}}
\end{vidu}
